\section{Motivation and Literature Position}
\label{sec:motivation}

Fast inverse calibration is useful when many option surfaces must be mapped to model
parameters repeatedly. Neural surrogates can reduce inference cost, but their value
cannot be assessed by speed alone. A valid comparison must preserve the same dynamics,
observation contract, constraints, split boundary, and metric families. It must also
distinguish structural validity from economic correctness and pricing fit from
parameter recovery.

The project's earlier development work sharpened these requirements. Official-NSE
screening rejected unsupported fixed maturities and extreme wings. Local Jacobian,
multi-start recovery, global-ambiguity, and complementary-observable studies showed
that low repricing error could coexist with materially displaced canonical vectors.
A subsequent R2/R3 experiment found that a third expiry improves conditioning on clean
synthetic data, but the predeclared stopping rule selected R2 because the clean
advantage disappeared at realistic noise levels. These studies are development or
historical context for the frozen benchmark; they do not establish unique
identification.

\subsection{Stochastic-volatility and affine pricing foundations}
\label{subsec:literature-stochastic-volatility}

Heston's stochastic-variance formulation provides a tractable alternative to a
constant-volatility Black--Scholes model by deriving closed-form option values with
applications across asset classes \citep{heston1993closed}. Affine transform methods
generalize this route to broader jump-diffusion classes and make factor-based
characteristic functions analytically useful \citep{duffie2000transform}. Stable
numerical evaluation of the square-root variance process remains important; the
Little-Heston-Trap representation is used for that reason in the production
implementation \citep{albrecher2007little}.

Richer volatility structure can improve observable fit in ways that one-factor models
cannot. Christoffersen, Heston, and Jacobs show that a two-factor specification can
represent independent movements in the level and slope of the index-option smirk and
report improved in-sample/out-of-sample fit relative to their benchmark
\citep{christoffersen2009shape}. Broader empirical comparisons motivate looking beyond
Black--Scholes without treating any single fitted model as economic truth
\citep{bakshi1997empirical}. Work comparing alternative stock-price dynamics likewise
separates model selection from parameter estimation \citep{chernov2003alternative}.

These observations support this paper's choice of a canonical two-factor Double Heston
model. The added factors are useful only if the inverse mapping remains scientifically
interpretable; improved observable fit alone does not establish recovery of ten latent
parameters.

\subsection{Calibration as an inverse problem}
\label{subsec:literature-calibration-identifiability}

Classical estimation treats option prices or integrated-variance summaries as data from
which latent dynamics must be inferred. Maximum likelihood has been applied to
stochastic-volatility diffusions \citep{aitsahalia2007maximum}, while conditional
moments of integrated volatility provide another estimator family
\citep{bollerslev2002estimating}. Such approaches address parameter inference but do not
remove finite-sample ambiguity: two parameterizations may imply nearly identical
observations while differing materially in individual components.

That distinction drives the evaluation design here. A low forward residual says that a
candidate reprices an observation vector under a chosen norm. It does not say that the
candidate equals the generating vector. The paper therefore reports range-scaled known-
truth RMSE separately from normalized repricing error, retains optimizer starts, and
conditions recovery statements on explicit tolerances.

\subsection{Neural acceleration and amortized calibration}
\label{subsec:literature-neural-calibration}

A recurring motivation for neural methods is computational cost. Liu et al.\ propose a
neural framework in which offline learning over model settings supports faster online
financial-model calibration, including high-dimensional stochastic-volatility cases
\citep{liu2019neural}. Bayer et al.\ learn a fast rough-stochastic-volatility price map
and then use conventional calibration in a second step \citep{bayer2019deep}. The
journal treatment of deep-learning volatility similarly couples learned pricing maps
with calibration \citep{horvath2021deep}.

Recent work refines these designs for speed, stability, and input representation. Sridi
and Bilokon apply differential machine learning to Heston calibration and study
regularization/generalization \citep{sridi2023applying}. Baschetti et al.\ combine
random grids with pointwise calibration for Heston and rough Bergomi, targeting grid
robustness without interpolation/extrapolation \citep{baschetti2024deep}. Biagini et
al.\ establish approximation rates for neural option-price maps, separating forward-map
approximation from finite-sample inversion \citep{biagini2024approximation}. Zhang et
al.\ learn prices and parameter sensitivities for gradient-based Heston calibration
\citep{zhang2025calibrating}, and Wu et al.\ preserve Fourier structure while using a
small network to accelerate a jump/rough-volatility calibration
\citep{wu2025efficient}. Methodological critiques also warn against equating neural fit
quality with valid calibrated behavior \citep{itkin2019deep}.

This literature motivates Model1 and Model2 but does not define their scientific claim.
Model1 is ordinary supervised inverse regression. Model2 adds hard constraints and a
differentiable Fourier-repricing loss; it remains repricing-informed, not PDE-informed.
The present synthetic protocol evaluates both against traditional calibration using
stored truth, frozen splits, failed-start retention, seed dispersion, runtime, and
conditioned recovery metrics.

\subsection{Physics-informed learning and option pricing}
\label{subsec:literature-pinn}

Equation-residual neural networks have a long history in ODE/PDE computation
\citep{lagaris1998artificial} and modern deep PDE solvers
\citep{sirignano2018dgm}. Physics-informed neural networks formalize the combination of
data fit and PDE residuals for forward and inverse problems
\citep{raissi2019physics}; reviews describe their uses and limitations across scientific
machine learning \citep{karniadakis2021physics,cuomo2022scientific}. Training failures
can arise even when a residual is mathematically well posed \citep{wang2022when}.

Finance applications include physics-informed Black--Scholes pricing
\citep{dhiman2023physics} and physics-informed volatility-surface prediction
\citep{kim2022physics}. Such work differs from Model2: a differentiable pricer or
repricing objective is not a Double Heston PDE residual. Genuine PDE-informed learning
must differentiate through state variables in the specified operator, verify terminal
and boundary conventions, and report stability and cost. This distinction is why Model3
is treated as separate pre-registered methodology rather than a renamed or tuned version
of Model2.

\subsection{Noise, incomplete observations, and evaluation scope}
\label{subsec:literature-noise-robustness}

Real surfaces are noisy, masked by activity/support rules, and economically redundant
across calls and puts once carry is fixed. Neural-calibration critiques therefore
emphasize accuracy and validity pitfalls rather than training loss alone
\citep{itkin2019deep}. Random-grid calibration explicitly targets interpolation and
extrapolation concerns \citep{baschetti2024deep}, while Bayesian neural-SDE calibration
uses posterior uncertainty and historical plus option information to express robustness
rather than a single point estimate \citep{cuchiero2025robust}.

These studies support controlled noise testing but do not supply this project's frozen
boundary. Completed evidence covers deterministic cohorts, strict zero-noise gates,
full-population neural evaluations through 1\% noise, and the clean traditional subset
gate. Positive-level traditional subset calibration remains pending; therefore no
three-way noise conclusion, OOD conclusion, G8 result, or Model3 result is made here.

\subsection{Position of this study}
\label{subsec:this-work-position}

Against this background, this work studies a narrowly scoped question: how do
traditional calibration, an ordinary ANN, and a constraint/repricing-informed inverse
network compare on a frozen two-factor Heston representation when generating vectors are
known? It separates forward-pricing correctness, observable fit, structural validity,
speed, seed/multi-start stability, and tolerance-conditioned parameter recovery. It does
not infer novelty from missing citations and does not claim that practical
non-identifiability is absent from prior literature. Instead, it reports where retained
ambiguity appears under a fixed DH contract and preserves a genuine PDE protocol as
future work.
